\documentclass[otchet]{../SCWorks}
% Тип обучения (одно из значений):
%    bachelor   - бакалавриат (по умолчанию)
%    spec       - специальность
%    master     - магистратура
% Форма обучения (одно из значений):
%    och        - очное (по умолчанию)
%    zaoch      - заочное
% Тип работы (одно из значений):
%    coursework - курсовая работа (по умолчанию)
%    referat    - реферат
%  * otchet     - универсальный отчет
%  * nirjournal - журнал НИР
%  * digital    - итоговая работа для цифровой кафедры
%    diploma    - дипломная работа
%    pract      - отчет о научно-исследовательской работе
%    autoref    - автореферат выпускной работы
%    assignment - задание на выпускную квалификационную работу
%    review     - отзыв руководителя
%    critique   - рецензия на выпускную работу

% * Добавлены вручную. За вопросами к @mchernigin
\usepackage{../preamble}

\begin{document}

% Кафедра (в родительном падеже)
\chair{математической кибернетики и компьютерных наук}

% Тема работы
\title{Анализ сложности быстрой и пирамидальной сортировок}

% Курс
\course{2}

% Группа
\group{251}

% Факультет (в родительном падеже) (по умолчанию "факультета КНиИТ")
\department{факультета компьютерных наук и информационных технологий}

% Специальность/направление код - наименование
% \napravlenie{02.03.02 "--- Фундаментальная информатика и информационные технологии}
% \napravlenie{02.03.01 "--- Математическое обеспечение и администрирование информационных систем}
% \napravlenie{09.03.01 "--- Информатика и вычислительная техника}
\napravlenie{09.03.04 "--- Программная инженерия}
% \napravlenie{10.05.01 "--- Компьютерная безопасность}

% Для студентки. Для работы студента следующая команда не нужна.
\studenttitle{студентки}

% Фамилия, имя, отчество в родительном падеже
\author{Потапкиной Маргариты Андреевны}

% Заведующий кафедрой 
\chtitle{доцент, к.\,ф.-м.\,н.}
\chname{С.\,В.\,Миронов}

% Руководитель ДПП ПП для цифровой кафедры (перекрывает заведующего кафедры)
% \chpretitle{
%     заведующий кафедрой математических основ информатики и олимпиадного\\
%     программирования на базе МАОУ <<Ф"=Т лицей №1>>
% }
% \chtitle{г. Саратов, к.\,ф.-м.\,н., доцент}
% \chname{Кондратова\, Ю.\,Н.}

% Научный руководитель (для реферата преподаватель проверяющий работу)
\satitle{старший преподаватель} %должность, степень, звание
\saname{М.\,И.\,Сафрончик}

% Руководитель практики от организации (руководитель для цифровой кафедры)
\patitle{доцент, к.\,ф.-м.\,н.}
\paname{И.\,А.\,Батраева}

% Руководитель НИР
\nirtitle{доцент, к.\,п.\,н.} % степень, звание
\nirname{И.\,А.\,Батраева}

% Семестр (только для практики, для остальных типов работ не используется)
\term{2}

% Наименование практики (только для практики, для остальных типов работ не
% используется)
\practtype{учебная}

% Продолжительность практики (количество недель) (только для практики, для
% остальных типов работ не используется)
\duration{2}

% Даты начала и окончания практики (только для практики, для остальных типов
% работ не используется)
\practStart{01.07.2024}
\practFinish{13.01.2024}

% Год выполнения отчета
\date{2026}

\maketitle

% Включение нумерации рисунков, формул и таблиц по разделам (по умолчанию -
% нумерация сквозная) (допускается оба вида нумерации)
\secNumbering

\tableofcontents

% Раздел "Обозначения и сокращения". Может отсутствовать в работе
% \abbreviations
% \begin{description}
%     \item ... "--- ...
%     \item ... "--- ...
% \end{description}

% Раздел "Определения". Может отсутствовать в работе
% \definitions

% Раздел "Определения, обозначения и сокращения". Может отсутствовать в работе.
% Если присутствует, то заменяет собой разделы "Обозначения и сокращения" и
% "Определения"
% \defabbr

\intro
В данном отчёте производится описание и анализ сложности быстрой и пирамидальной сортировок.

\section{Быстрая сортировка}
\small
\inputminted[breaklines]{cpp}{./quick_sort.cpp}
\normalsize

В быстрой сортировке происходит выбор опорного элемента, затем элементы массива переупорядочиваются так, что меньшие опорного оказываются слева, а большие "--- справа. Этот процесс выполняется за время, пропорциональное размеру текущего массива. После этого функция рекурсивно вызывается от обеих половин.

Функция трудоёмкости быстрой сортировки является порядково"=зависимой: в лучшем случае, при оптимальном выборе опорного элемента, массив каждый раз будет делиться примерно пополам. Тогда можно построить следующее рекуррентное соотношение:
\begin{equation}
	t(n) =
\begin{cases}
	c, \text{если } n = 1; \\
	2 \cdot t(n / 2) + bn, \text{если } n > 1.
\end{cases}
\end{equation}

Согласно следствию из теоремы 1, в таком случае алгоритм будет иметь асимптотику $O(n \cdot log(n))$.

Однако существуют случаи, когда быстрая сортировка может работать за время $O(n^2)$. Если в качестве опорного элемента выбирается первый/последний элемент массива, то к худшему случаю приводит массив, отсортированный в прямом или обратном порядке. И для случая, когда мы выбираем средний элемент в качестве опорного, можно подобрать контртест:
\begin{minted}{cpp}
for (int i = 0; i < n; i++) swap(a[i], a[i / 2]);
\end{minted}

В данном случае роль опорного элемента всегда будет играть максимальный, а значит, массив каждый раз будет разделяться на части размера $1$ и $n - 1$. Тогда рекуррентное соотношение будет иметь вид:

\begin{equation}
	t(n) =
\begin{cases}
	c, \text{если } n = 1; \\
	t(n - 1) + bn, \text{если } n > 1.
\end{cases}
\end{equation}

Согласно следствию из теоремы 2, в данном случае алгоритм будет иметь асимптотику $O(n^2)$.

Однако в моей реализации быстрой сортировки опорный элемент выбирается случайным образом. Это усложняет подбор теста, который приведёт к квадратичному времени работы, и можно сказать, что в среднем такая реализация работает за $O(n \cdot log(n))$.

Здесь $n$ "--- количество элементов в исходном массиве.

\section{Пирамидальная сортировка}
\small
\inputminted[breaklines]{cpp}{./heap_sort.cpp}
\normalsize

Работа функции \textit{heapify} основана на принципе <<разделяй и властвуй>>. Дети $i$"=го ребёнка в двоичной куче, которая в данном случае хранится в виде массива, имеют номера $2i + 1$ и $2i + 2$. В результате при каждом вызове значение $i$ увеличивается хотя бы в $2$ раза, при этом оно не может стать больше, чем $n$, таким образом будет сделано $log(n)$ вызовов. Один вызов выполняется за $O(1)$, а значит, время работы функции \textit{heapify} составляет $O(log(n))$.

В пирамидальной сортировке предварительно функция просеивания вызывается $n$ раз, асимптотика этих вызовов составляет $n \cdot log(n)$ по правилу произведения. Затем ещё $n$ раз происходит просеивание, каждый раз в результате него максимальный элемент поднимается наверх (оказывается на первом месте), и за $O(1)$ происходит его перемещение на соответствующую позицию (в конец необработанной части массива). Итого эта часть кода работает за $O(n \cdot log(n))$.

\textbf{Итого асимптотику алгоритма можно оценить как $O(n \cdot log(n))$.}

Здесь $n$ "--- количество элементов в исходном массиве.

% После введения — серии \section, \subsection и т.д.

% \conclusion

% Отобразить все источники. Даже те, на которые нет ссылок.
% \nocite{*}

\bibliographystyle{ugost2008}
\bibliography{thesis}

% Окончание основного документа и начало приложений Каждая последующая секция
% документа будет являться приложением
\appendix

\end{document}
